\documentclass[tikz,border=8pt]{standalone}
\usepackage{ctex}
\usepackage{amsmath}
\usetikzlibrary{arrows.meta, shapes.geometric, shapes.misc,
                positioning, fit, backgrounds, calc}

\begin{document}
\begin{tikzpicture}[
  node distance=1.4cm and 2.0cm,
  box/.style={
    rectangle, rounded corners=4pt,
    draw=black!70, fill=blue!8,
    minimum width=2.6cm, minimum height=0.9cm,
    font=\small, align=center,
  },
  formula/.style={
    rectangle, rounded corners=2pt,
    draw=gray!50, fill=yellow!10,
    minimum width=4.2cm, minimum height=0.75cm,
    font=\footnotesize, align=center,
  },
  arrow/.style={-{Stealth[scale=1.1]}, thick, gray!80},
  label/.style={font=\scriptsize, gray!70},
]

% ====== 节点 ======
\node[box, fill=green!15, draw=green!60!black] (g)
  {随机梯度 $g_t$};

\node[box, right=2.5cm of g] (m)
  {一阶矩 $m_t$};

\node[box, below=0.9cm of m] (v)
  {二阶矩 $v_t$};

\node[box, right=2.5cm of m] (mhat)
  {偏差修正\\$\hat{m}_t$};

\node[box, below=0.9cm of mhat] (vhat)
  {偏差修正\\$\hat{v}_t$};

\node[box, right=2.5cm of mhat, fill=orange!15, draw=orange!60!black,
      minimum height=1.5cm] (delta)
  {参数更新\\$\Delta\theta_t$};

\node[box, right=1.8cm of delta, fill=red!15, draw=red!60!black] (theta)
  {$\theta_{t+1}$};

% ====== 公式注释（formula 节点）======
\node[formula, above=0.5cm of m] (fm)
  {$m_t = \beta_1 m_{t-1} + (1-\beta_1)g_t$};

\node[formula, below=0.5cm of v] (fv)
  {$v_t = \beta_2 v_{t-1} + (1-\beta_2)g_t^2$};

\node[formula, above=0.5cm of mhat] (fmh)
  {$\hat{m}_t = \dfrac{m_t}{1-\beta_1^t}$};

\node[formula, below=0.5cm of vhat] (fvh)
  {$\hat{v}_t = \dfrac{v_t}{1-\beta_2^t}$};

\node[formula, above=0.65cm of delta] (fdelta)
  {$\Delta\theta_t = -\dfrac{\alpha\,\hat{m}_t}{\sqrt{\hat{v}_t}+\epsilon}$};

\node[formula, below=0.65cm of delta] (fadamw)
  {(AdamW: $-\alpha\lambda\theta_t$ 直接加)};

% ====== 箭头 ======
% g -> m 和 g -> v
\draw[arrow] (g.east) -- (m.west)
  node[label, midway, above] {EMA};

\draw[arrow] (g.east) -- ++(1.0,0) |- (v.west)
  node[label, pos=0.3, right] {EMA};

% m -> mhat, v -> vhat
\draw[arrow] (m.east) -- (mhat.west)
  node[label, midway, above] {修正};

\draw[arrow] (v.east) -- (vhat.west)
  node[label, midway, above] {修正};

% mhat -> delta, vhat -> delta
\draw[arrow] (mhat.east) -- (delta.west |- mhat)
  node[label, midway, above] {分子};

\draw[arrow] (vhat.east) -- (delta.west |- vhat)
  node[label, midway, below] {分母};

% delta -> theta
\draw[arrow] (delta.east) -- (theta.west)
  node[label, midway, above] {更新};

% 公式连线（虚线）
\draw[gray!40, dashed, thin] (fm.south) -- (m.north);
\draw[gray!40, dashed, thin] (fv.north) -- (v.south);
\draw[gray!40, dashed, thin] (fmh.south) -- (mhat.north);
\draw[gray!40, dashed, thin] (fvh.north) -- (vhat.south);
\draw[gray!40, dashed, thin] (fdelta.south) -- (delta.north);
\draw[gray!40, dashed, thin] (fadamw.north) -- (delta.south);

% ====== 标题 ======
\node[font=\normalsize, above=1.8cm of m, xshift=3cm]
  {Adam 更新公式分解流程图};

% ====== 超参数注释 ======
\node[font=\scriptsize, gray, below=0.2cm of fadamw, xshift=0cm]
  {默认：$\beta_1{=}0.9,\ \beta_2{=}0.999,\ \epsilon{=}10^{-8},\ \alpha{=}10^{-3}$};

\end{tikzpicture}
\end{document}
